\documentclass[12pt,a4paper]{article}

\usepackage{template}
\usepackage{amsmath}
\usepackage{amssymb}
\usepackage{booktabs}
\usepackage{fontspec}
\setmainfont{PT Serif}
\setsansfont{PT Sans}
\setmonofont{PT Mono}

\hypersetup{
  colorlinks=true,
  linkcolor=blue!45!black,
  urlcolor=blue!45!black,
  citecolor=blue!45!black,
  pdfauthor={Иван Афанасьев},
  pdftitle={Задача 3. Немного о перестановках}
}

\begin{document}
\sloppy

\begin{center}
  {\LARGE\bfseries Задача 3. Немного о перестановках}\\[4pt]
  {\large Иван Афанасьев}\\[2pt]
\end{center}

\vspace{6pt}
\begin{center}
\fbox{\begin{minipage}{0.92\textwidth}
\small
\textbf{Использование ИИ.}
Языковые модели (GLM-5.2, OpenCode-CI-Agent) применялись для
независимой верификации ответов (через полный перебор всех
циклов для $n\le12$), рецензии кода (\texttt{p3\_brute.cpp},
\texttt{p3\_construct.cpp}) и оформления \LaTeX-документа.
\end{minipage}}
\end{center}
\vspace{10pt}

\section{Постановка задачи}

Дано $n\ge2$. Перестановку $\sigma$ чисел $1,\dots,n$ применяем многократно;
$\sigma^{k}$~-- результат $k$ применений. $T_n$~-- множество перестановок,
у которых для каждого $i=1,\dots,n-1$ найдётся натуральное $k$ (своё для
каждого $i$) с $\sigma^{k}(i)=i+1$. Функция $f(\tau)$ считает, сколько
элементов $i\in\{1,\dots,n-1\}$ переходят в $i+1$ сразу: $\tau(i)=i+1$.
Для каждого $n$ нужно найти:
\begin{enumerate}[label=\arabic*.]
  \item $\max\limits_{\sigma\in T_n}\ \min\{f(\sigma^{k})\mid k\in\mathbb N\}$;
  \item $\max\limits_{\sigma\in T_n}\ \min\{f(\sigma^{2k-1})\mid k\in\mathbb N\}$;
  \item $\min\limits_{\sigma\in T_n}\ \max\{f(\sigma^{2k})\mid k\in\mathbb N\}$;
  \item $\min\limits_{\sigma\in T_n}\ \max\{f(\sigma^{2k-1})\mid k\in\mathbb N\}$;
  \item $\min\limits_{\sigma\in T_n}\ \max\{f(\sigma^{k})\mid k\in\mathbb N\}$.
\end{enumerate}

Все значения подтверждены полным перебором для $n=2,\dots,12$
(\texttt{p3\_brute.cpp}) и явной конструкцией для произвольного $n$
(\texttt{p3\_construct.cpp}).

\section{Шаг 1. Что за перестановки лежат в \texorpdfstring{$T_n$}{T\_n}}

\textbf{Утверждение.} $T_n$~-- это в точности перестановки, являющиеся
одним циклом длины $n$ (все $n$ элементов образуют один <<хоровод>>:
$\sigma$ переводит каждого в следующего по кругу).

\emph{Доказательство.} Многократное применение $\sigma$ гоняет элемент $i$
по его циклу, поэтому <<$\sigma^{k}(i)=i+1$ для некоторого $k$>> означает
ровно то, что $i+1$ стоит в том же цикле, что и $i$. Если это верно для
всех $i=1,\dots,n-1$, то $1$ и $2$ в одном цикле, $2$ и $3$ в одном
цикле, \dots~-- все $n$ элементов в одном цикле, т.\,е. $\sigma$~--
один цикл длины $n$. Обратно, если $\sigma$~-- один цикл, то из любого
элемента многократным применением можно попасть в любой другой.
\hfill$\blacksquare$

\section{Шаг 2. Модель: шаги по кругу и остатки}

Пусть $\sigma\in T_n$~-- цикл. Расставим числа $1,\dots,n$ по кругу в том
порядке, в котором их обходит $\sigma$, и пронумеруем места $0,1,\dots,n-1$
(место числа $1$~-- нулевое). Пусть число $i$ стоит на месте $L(i)$.
Применение $\sigma$~-- это сдвиг на одно место по кругу, применение
$\sigma^{k}$~-- сдвиг на $k$ мест. Поэтому
\[
  \sigma^{k}(i)=i+1
  \iff
  k\equiv L(i+1)-L(i)\pmod n .
\]
Обозначим $d_i$~-- остаток от деления $L(i+1)-L(i)$ на $n$; так как места
различны, $d_i\in\{1,\dots,n-1\}$. Получаем главный вывод:
\[
  f(\sigma^{k})=\#\{\,i:\ d_i=r\,\},\qquad\text{где $r$~-- остаток $k$ по модулю $n$};
\]
обозначим эту величину $m_r$ (при $r=0$ получается тождественная
перестановка и $f=0$). Так как чисел $d_1,\dots,d_{n-1}$ ровно $n-1$,
\[
  m_1+m_2+\dots+m_{n-1}=n-1. \tag{$\ast$}
\]

\emph{Какие наборы $d_i$ возможны?} Суммы $S_0=0$, $S_i=d_1+\dots+d_i$~--
это в точности места $L(1),L(2),\dots,L(n)$. Значит набор
$(d_1,\dots,d_{n-1})$ отвечает какому-то циклу тогда и только тогда, когда
остатки чисел $S_0,S_1,\dots,S_{n-1}$ по модулю $n$ попарно различны
(тогда они занимают все $n$ мест). Выбор $\sigma\in T_n$~-- это выбор
такого набора.

\emph{Какие остатки даёт чётность $k$?} Остаток $k$ по модулю $n$ при
нечётных/чётных $k$:
\begin{itemize}
  \item $n$ нечётно: и нечётные, и чётные $k$ дают \emph{все} остатки
        $0,\dots,n-1$ (прибавление нечётного $n$ меняет чётность $k$, не
        меняя остатка);
  \item $n$ чётно: нечётные $k$ дают ровно нечётные остатки, чётные
        $k$~-- ровно чётные (включая $0$ при $k=n$).
\end{itemize}

Теперь все пять вопросов~-- это вопросы о частотах $m_r$.

\section{Шаг 3. Ответы}

Обозначим $M=\max_r m_r$; из $(\ast)$ видно $M\ge1$.

\paragraph{Вопрос 1: ответ $0$.} При $k=n$ получаем тождественную
перестановку и $f(\sigma^{n})=0$, поэтому минимум по всем $k$ равен $0$
для \emph{любой} $\sigma$. \hfill$\blacksquare$

\paragraph{Вопрос 5: ответ $1$ при чётном $n$, $2$ при нечётном.}
Максимум по всем $k$ равен $M$, ищем наименьшее возможное $M$.
\begin{itemize}
  \item \emph{$M=1$ при нечётном $n$ невозможно:} тогда каждое из чисел
        $1,\dots,n-1$ встречается среди $d_i$ ровно один раз, и
        \[
          S_{n-1}=1+2+\dots+(n-1)=\frac{n(n-1)}2 .
        \]
        При нечётном $n$ число $\frac{n-1}2$ целое, значит $S_{n-1}$
        делится на $n$ и даёт тот же остаток, что $S_0=0$,~-- а остатки
        обязаны быть различны. Значит $M\ge2$.
  \item \emph{Достижимость} ($M=1$ при чётном $n$, $M=2$ при нечётном)~--
        зигзаг-конструкция в Шаге~4.
\end{itemize}
\hfill$\blacksquare$

\paragraph{Вопросы 3 и 4 при нечётном $n$: оба ответа $2$.} При нечётном
$n$ и чётные, и нечётные $k$ пробегают все остатки, поэтому оба максимума
равны $M$, а наименьшее $M$ равно $2$ (вопрос~5). \hfill$\blacksquare$

\paragraph{Вопрос 3 при чётном $n$: ответ $0$.} Чётные $k$ дают только
чётные остатки. Возьмём простейший цикл $\sigma=(1\,2\,\dots\,n)$
(каждый переходит в следующего): у него все $d_i=1$, т.\,е. $m_1=n-1$ и
$m_r=0$ при $r\ge2$. Все чётные остатки имеют частоту $0$, значит максимум
по чётным $k$ равен $0$~-- меньше не бывает. \hfill$\blacksquare$

\paragraph{Вопрос 4 при чётном $n$: ответ $1$.} Нечётные $k$ дают нечётные
остатки; ищем наименьший возможный $\max_{r\text{ нечёт}}m_r$.
\begin{itemize}
  \item \emph{Меньше $1$ нельзя:} если бы все $d_i$ были чётными, то и все
        суммы $S_i$ были бы чётными; но чётных остатков по модулю чётного
        $n$ только $n/2$, а различных сумм нужно $n$. Значит какой-то
        нечётный остаток встречается, и максимум по нечётным $r$ не
        меньше $1$.
  \item \emph{Ровно $1$}~-- у зигзага (Шаг~4). \hfill$\blacksquare$
\end{itemize}

\paragraph{Вопрос 2: ответ $1$ при чётном $n$, $0$ при нечётном.}
\begin{itemize}
  \item \emph{$n$ нечётно:} среди нечётных $k$ есть дающие остаток $0$
        (например $k=n$), для них $f=0$; минимум равен $0$ у любой
        $\sigma$.
  \item \emph{$n$ чётно:} минимум по нечётным $k$~-- это
        $\min_{r\text{ нечёт}}m_r$. Если бы он был $\ge2$, то сумма
        частот только по $n/2$ нечётным остаткам была бы $\ge n$,
        вопреки $(\ast)$ (всего $n-1$). Значит не больше $1$; ровно
        $1$~-- у зигзага, где каждый нечётный остаток встречается один
        раз. \hfill$\blacksquare$
\end{itemize}

\section{Шаг 4. Зигзаг: явная наилучшая перестановка}

Выберем места $S_0,S_1,S_2,\dots$ <<зигзагом>>:
\[
  0,\;1,\;-1,\;2,\;-2,\;3,\;-3,\dots\pmod n .
\]
Эти остатки попарно различны: при чётном $n$ это
$\{0,\pm1,\dots,\pm(\tfrac n2-1),\tfrac n2\}$, при нечётном~--
$\{0,\pm1,\dots,\pm\tfrac{n-1}2\}$. Разности равны
\[
  d_1=1,\ d_2=-2,\ d_3=3,\ d_4=-4,\dots,\qquad d_i=(-1)^{i+1}i .
\]
\begin{itemize}
  \item \textbf{$n$ чётно:} по модулю $n$ нечётные $i$ дают нечётные
        остатки $1,3,\dots,n-1$, а чётные $i$~-- остатки $n-i$, т.\,е.
        все чётные $2,4,\dots,n-2$. Каждый ненулевой остаток встречается
        ровно один раз: все $m_r=1$, $M=1$. Одна эта перестановка даёт
        оптимум сразу в вопросах 2, 4 и 5.
  \item \textbf{$n$ нечётно:} нечётные $i$ дают остатки $1,3,\dots,n-2$,
        а чётные $i$~-- остатки $n-i$, и при нечётном $n$ это \emph{те
        же} нечётные числа $n-2,n-4,\dots,1$. Каждый нечётный остаток
        встречается ровно дважды, чётные~-- ни разу: $M=2$. Эта
        перестановка даёт оптимум в вопросах 3, 4 и 5.
\end{itemize}
Вместе с простейшим циклом из вопроса~3 (чётное $n$) этим закрыты все
случаи достижимости.

Например, при $n=4$ зигзагу отвечают места $0,1,3,2$ для чисел $1,2,3,4$,
т.\,е. цикл $1\to2\to4\to3\to1$: у него $f(\sigma)=1$ ($1\to2$),
$f(\sigma^{2})=1$ ($2\to3$), $f(\sigma^{3})=1$ ($3\to4$),
$f(\sigma^{4})=0$.

\section{Вычислительная проверка}

\texttt{p3\_brute.cpp} перебирает \emph{все} циклы длины $n$ (их
$(n-1)!$), честно вычисляет все пять величин и для $n=2,\dots,12$
воспроизводит ответы на все пять вопросов без единого расхождения.
\texttt{p3\_construct.cpp} строит зигзаг для каждого $n$ от $2$ до $40$,
проверяет различность мест, пересчитывает частоты $m_r$ напрямую через
степени перестановки и подтверждает $M=1$ (чётные $n$) и $M=2$ (нечётные).

\begin{verbatim}
g++ -O2 -std=c++17 -o p3_brute p3_brute.cpp && ./p3_brute 12
g++ -O2 -std=c++17 -o p3_construct p3_construct.cpp && ./p3_construct 2 40
\end{verbatim}

\emph{Замечание для интересующихся.} Наборы разностей, в которых каждый
ненулевой остаток встречается ровно один раз, в комбинаторике называются
\emph{секвенцированиями}; то, что для чисел по модулю $n$ они существуют
ровно при чётных $n$,~-- известная теорема (B.~Gordon, 1961). Обе нужные
нам половины этой теоремы полностью доказаны выше школьными средствами.

\end{document}
